\subsubsection{Sous-réseau invariant et unité quadratique de l'opérateur agrégé}
\label{mg02:subreticula}

L'agrégation par classes de résidus conserve une structure entière qui
précise la décomposition spectrale précédente. L'opérateur de départ est
\(M_\Pi\), obtenu en moyennant le feuillet multiplicatif d'APP sur
les blocs \(C_a\times C_b\). Son domaine entier est \(\mathbb Z^3\),
avec les coordonnées ordonnées par \(C_1,C_2,C_0\).

\begin{proposition}[Décomposition sur un sous-réseau d'indice deux]
Soient
\[
 v_-=(1,-1,0)^{\mathsf T},\qquad
 v_+=(1,1,0)^{\mathsf T},\qquad v_0=(0,0,1)^{\mathsf T},
\]
et soit \(L=\mathbb Zv_-\oplus\mathbb Zv_+\oplus\mathbb Zv_0\).
Alors
\begin{equation}
 L=\{(x,y,z)\in\mathbb Z^3:x\equiv y\pmod2\},
 \qquad [\mathbb Z^3:L]=2.
 \label{mg02:indice}
\end{equation}
Le sous-réseau \(L\) est invariant par \(M_\Pi\). Dans la base
ordonnée \((v_-,v_+,v_0)\), la restriction a pour matrice
\begin{equation}
 [M_\Pi|_L]=(-1)\oplus3H,\qquad
 H=\begin{pmatrix}3&2\\4&3\end{pmatrix},\qquad \det H=1.
 \label{mg02:restriccion}
\end{equation}
\end{proposition}

\begin{proof}
La combinaison \(av_-+bv_++cv_0=(a+b,-a+b,c)\) possède ses deux
premières coordonnées de même parité. Réciproquement, pour un vecteur
possédant cette parité, les coordonnées dans la base indiquée sont
\(a=(x-y)/2\), \(b=(x+y)/2\) et \(c=z\), toutes entières. La matrice
de changement de base et son inverse rationnelle sont
\[
 P=\begin{pmatrix}1&1&0\\-1&1&0\\0&0&1\end{pmatrix},\qquad
 P^{-1}=\begin{pmatrix}1/2&-1/2&0\\1/2&1/2&0\\0&0&1\end{pmatrix}.
\]
Dans cet ordre des colonnes, \(\det P=2\), ce qui donne l'indice de
\eqref{mg02:indice}. L'action sur les trois générateurs est
\[
 M_\Pi v_-=-v_-,\qquad
 M_\Pi v_+=9v_++12v_0,\qquad
 M_\Pi v_0=6v_++9v_0.
\]
Ces égalités prouvent à la fois l'invariance et l'identité
\[
 P^{-1}M_\Pi P=
 \begin{pmatrix}-1&0&0\\0&9&6\\0&12&9\end{pmatrix}
 =(-1)\oplus3H.
\]
Enfin, \(\det H=9-8=1\).
\end{proof}

\begin{corollary}[Itération de Pell dans le sous-module de rang deux]
Pour \(n\geq0\), il existe des entiers \(a_n,b_n\) tels que
\begin{equation}
 H^n=\begin{pmatrix}a_n&b_n\\2b_n&a_n\end{pmatrix},\qquad
 a_n+b_n\sqrt2=(3+2\sqrt2)^n,
 \qquad a_n^2-2b_n^2=1.
 \label{mg02:pell}
\end{equation}
En particulier, \((a_0,b_0)=(1,0)\), \((a_1,b_1)=(3,2)\) et
\((a_2,b_2)=(17,12)\). Les valeurs propres de \(H\) sont les unités
quadratiques conjuguées \(3\pm2\sqrt2\).
\end{corollary}

\begin{proof}
La multiplication des matrices de la forme indiquée correspond à
\[
 (a,b)(c,d)=(ac+2bd,ad+bc),
\]
qui est la loi de multiplication de \(a+b\sqrt2\). Ainsi, les puissances
de \(H\) satisfont
\(a_{n+1}=3a_n+4b_n\), \(b_{n+1}=2a_n+3b_n\), avec les données
initiales précédentes. La norme multiplicative de \(3+2\sqrt2\) vaut
un ; par conséquent, \(a_n^2-2b_n^2=1\) à tout degré. Le polynôme
caractéristique de \(H\), \(\lambda^2-6\lambda+1\), donne les deux
unités conjuguées. L'identité
\(H^{\mathsf T}\operatorname{diag}(2,-1)H=
\operatorname{diag}(2,-1)\) exprime la même conservation sous la forme
d'une forme bilinéaire entière.
\end{proof}

La séparation entière réside donc dans \(L\) ; la même matrice \(P\) donne une décomposition complète sur \(\mathbb R^3\). Le facteur \(3\) de la restriction conserve l'échelle de la moyenne APP : son itération de rang deux est \(3^nH^n\). Le bloc central \(B_c=\left(\begin{smallmatrix}7&2\\2&7\end{smallmatrix}\right)\)
réside en revanche dans deux positions adjacentes du tableau et satisfait
\(B_c^{\mathsf T}\operatorname{diag}(1,-1)B_c=
45\operatorname{diag}(1,-1)\). La compression en trois classes et la
restriction centrale à deux positions conservent ainsi leurs domaines,
bases et normalisations respectifs. La structure de Pell procède de
la première opération au moyen de \eqref{mg02:restriccion}.
